\hypertarget{numeric-mathematical-and-cryptographic-randomness}{%
\chapter{Numeric, Mathematical, and Cryptographic
Randomness}\label{numeric-mathematical-and-cryptographic-randomness}}

Python provides a robust suite of modules for numerical computing,
ranging from standard floating-point math to arbitrary-precision
decimals and cryptographically secure random number generation.

\hypertarget{decimal-control-over-precision}{%
\subsection{\texorpdfstring{35.1 \texttt{decimal}: Control over
Precision}{35.1 decimal: Control over Precision}}\label{decimal-control-over-precision}}

The standard \passthrough{\lstinline!float!} in Python is a 64-bit IEEE
754 double, which suffers from precision issues (e.g.,
\passthrough{\lstinline"0.1 + 0.2 != 0.3"}). The
\passthrough{\lstinline!decimal!} module provides a
\passthrough{\lstinline!Decimal!} type for correctly-rounded decimal
floating-point arithmetic.

\hypertarget{the-decnumber-c-library}{%
\subsubsection{\texorpdfstring{1. The \texttt{decNumber} C
Library}{1. The decNumber C Library}}\label{the-decnumber-c-library}}

In CPython, the \passthrough{\lstinline!decimal!} module is implemented
as \passthrough{\lstinline!\_decimal!}, which is a wrapper around the
\textbf{decNumber} library. This allows for extremely fast decimal
arithmetic that follows the General Decimal Arithmetic Specification.

\hypertarget{contexts-and-precision}{%
\subsubsection{2. Contexts and Precision}\label{contexts-and-precision}}

You can control the global or local precision using
\passthrough{\lstinline!getcontext()!}:

\begin{lstlisting}[language=Python]
from decimal import Decimal, getcontext
getcontext().prec = 50  # 50 digits of precision
print(Decimal(1) / Decimal(7))
\end{lstlisting}

\begin{itemize}
\tightlist
\item
  \textbf{Performance Note}: While \passthrough{\lstinline!\_decimal!}
  is fast, it is still significantly slower than hardware-native
  \passthrough{\lstinline!float!}. Use it for financial applications or
  cases where exact decimal representation is mandatory.
\end{itemize}

\hypertarget{fractions-exact-rational-numbers}{%
\subsection{\texorpdfstring{35.2 \texttt{fractions}: Exact Rational
Numbers}{35.2 fractions: Exact Rational Numbers}}\label{fractions-exact-rational-numbers}}

The \passthrough{\lstinline!fractions!} module provides support for
rational number arithmetic. * \textbf{Internals}: A
\passthrough{\lstinline!Fraction!} object stores two integers: a
numerator and a denominator. It automatically reduces the fraction to
its lowest terms using the Greatest Common Divisor (GCD). *
\textbf{Exactness}: Unlike \passthrough{\lstinline!float!} or
\passthrough{\lstinline!decimal!}, \passthrough{\lstinline!Fraction!}
can represent \passthrough{\lstinline!1/3!} exactly without any rounding
error.

\hypertarget{math-and-cmath-the-c-standard-library-wrappers}{%
\subsection{\texorpdfstring{35.3 \texttt{math} and \texttt{cmath}: The C
Standard Library
Wrappers}{35.3 math and cmath: The C Standard Library Wrappers}}\label{math-and-cmath-the-c-standard-library-wrappers}}

\begin{itemize}
\tightlist
\item
  \textbf{\passthrough{\lstinline!math!}}: Provides access to the
  mathematical functions defined by the C standard for real numbers
  (\passthrough{\lstinline!sin!}, \passthrough{\lstinline!cos!},
  \passthrough{\lstinline!log!}, \passthrough{\lstinline!sqrt!}).
\item
  \textbf{\passthrough{\lstinline!cmath!}}: Provides the same functions
  for complex numbers.
\item
  \textbf{Optimization}: These functions are thin wrappers around the
  host C library. They are highly optimized and release the GIL for
  heavy calculations (though most are too fast for the release overhead
  to be worth it).
\end{itemize}

\hypertarget{random-pseudorandom-number-generation}{%
\subsection{\texorpdfstring{35.4 \texttt{random}: Pseudorandom Number
Generation}{35.4 random: Pseudorandom Number Generation}}\label{random-pseudorandom-number-generation}}

The \passthrough{\lstinline!random!} module is a \textbf{Pseudorandom
Number Generator (PRNG)}. It is deterministic if you know the seed.

\hypertarget{the-mersenne-twister-mt19937}{%
\subsubsection{1. The Mersenne Twister
(MT19937)}\label{the-mersenne-twister-mt19937}}

Historically, Python used the Mersenne Twister as its primary PRNG. *
\textbf{Period}: \(2^{19937} - 1\). * \textbf{State}: It maintains a
large state (624 integers). * \textbf{Weakness}: It is not
cryptographically secure; observing a sufficient number of outputs
allows an attacker to predict future values.

\hypertarget{pcg64-python-3.13}{%
\subsubsection{2. PCG64 (Python 3.13+)}\label{pcg64-python-3.13}}

Modern Python versions have introduced more modern PRNGs like PCG64,
which offer better statistical properties and smaller state.

\hypertarget{secrets-cryptographic-security}{%
\subsection{\texorpdfstring{35.5 \texttt{secrets}: Cryptographic
Security}{35.5 secrets: Cryptographic Security}}\label{secrets-cryptographic-security}}

For security-sensitive applications (passwords, tokens), you must use
the \passthrough{\lstinline!secrets!} module. * \textbf{Internals}:
\passthrough{\lstinline!secrets!} uses the OS's cryptographically secure
source of randomness (\passthrough{\lstinline!/dev/urandom!} on Unix,
\passthrough{\lstinline!CryptGenRandom!} on Windows). * \textbf{Why?}:
Unlike \passthrough{\lstinline!random!}, the output of
\passthrough{\lstinline!secrets!} is not predictable even if an attacker
sees millions of previous values.

\begin{center}\rule{0.5\linewidth}{0.5pt}\end{center}

\begin{center}\rule{0.5\linewidth}{0.5pt}\end{center}
